% EN4554 Assignment 2 - Report
% Compile with: pdflatex report.tex  (run twice for references)
\documentclass[11pt,a4paper]{article}

\usepackage[T1]{fontenc}
\usepackage[utf8]{inputenc}
\usepackage{lmodern}
\usepackage[margin=2.5cm]{geometry}
\usepackage{graphicx}
\usepackage{booktabs}
\usepackage{amsmath}
\usepackage{amssymb}
\usepackage{xcolor}
\usepackage{listings}
\usepackage{caption}
\usepackage{float}
\usepackage[hidelinks]{hyperref}

% ---------------------------------------------------------------- figures
% Figures live in figures/ as PDF. If a file is missing, a labelled grey box
% is drawn instead, so the report always compiles.
\graphicspath{{figures/}}

\newcommand{\figbox}[1]{%
  \fcolorbox{gray!50}{gray!8}{%
    \begin{minipage}[c][4.5cm][c]{0.92\linewidth}
      \centering\small\texttt{\detokenize{figures/#1.pdf}}\\[2pt]
      \itshape(figure not found -- place the file here)
    \end{minipage}}%
}

% \fig{filename}{width fraction}{caption}
\newcommand{\fig}[3]{%
  \begin{figure}[H]
    \centering
    \IfFileExists{figures/#1.pdf}{\includegraphics[width=#2\linewidth]{#1}}{\figbox{#1}}
    \caption{#3}
    \label{fig:#1}
  \end{figure}%
}

% ---------------------------------------------------------------- code style
\definecolor{codebg}{RGB}{247,249,252}
\definecolor{codekw}{RGB}{164,36,154}
\definecolor{codestr}{RGB}{26,127,55}
\definecolor{codecom}{RGB}{107,122,153}

\lstdefinestyle{py}{
  language=Python,
  backgroundcolor=\color{codebg},
  basicstyle=\ttfamily\scriptsize,
  keywordstyle=\color{codekw},
  stringstyle=\color{codestr},
  commentstyle=\itshape\color{codecom},
  numbers=none,
  frame=leftline,
  framerule=1.2pt,
  rulecolor=\color{gray!60},
  framesep=6pt,
  xleftmargin=8pt,
  breaklines=true,
  showstringspaces=false,
  captionpos=b,
  aboveskip=1em,
  belowskip=1em,
}
\lstset{style=py}

\captionsetup{font=small}
\setlength{\parskip}{0.5em}
\setlength{\parindent}{0pt}

% ---------------------------------------------------------------- title
\title{\vspace{-1.5cm}\textbf{EN4554 --- Assignment 2 on CNNs}}
\author{Lakshan R.G.R \\ Index: 220355M}
\date{\today}

\begin{document}
\maketitle

\begin{center}
\small
GitHub: \url{https://github.com/ravindu374/EN4554_Assignment_2}
\end{center}

\vspace{-0.5em}
\hrule
\vspace{1em}

% =====================================================================
\section{Setup}

All experiments use CIFAR-10. The 50{,}000 training images are split once into
45{,}000 for training and 5{,}000 for validation using a fixed permutation, and
this split is reused everywhere so results stay comparable. Images are
normalised per channel,
%
\begin{equation}
x' = \frac{x - \mu_c}{\sigma_c},
\end{equation}
%
with $\mu_c$ and $\sigma_c$ computed from the training split only, so no
information leaks out of validation or test. Measured values were
$\mu = (0.4912, 0.4820, 0.4464)$ and $\sigma = (0.2472, 0.2436, 0.2617)$. All
random seeds are fixed. The test set is read exactly once, in Task 4.

The whole dataset is held on the GPU as \texttt{uint8} (138\,MB) and normalised
per batch, which removes \texttt{DataLoader} worker overhead. Every experiment
writes its history to \texttt{results/*.json}, so plots can be regenerated
without retraining.

% =====================================================================
\section{Task 1: Backpropagation}

\subsection{Implementation}

Each layer stores what its backward pass needs during the forward pass. For a
linear layer $y = xW + b$ with incoming gradient $g = \partial L/\partial y$:
%
\begin{equation}
\frac{\partial L}{\partial x} = g W^{\!\top}, \qquad
\frac{\partial L}{\partial W} = x^{\!\top} g, \qquad
\frac{\partial L}{\partial b} = \sum_{n=1}^{N} g_n .
\end{equation}
%
ReLU acts as a gate, passing gradient only where it passed signal:
$\partial L/\partial x = g \odot \mathbb{1}[x>0]$.

Softmax and cross-entropy are implemented as a single fused block. For logits
$z \in \mathbb{R}^{N \times C}$ and labels $y$,
%
\begin{equation}
L = \frac{1}{N}\sum_{n}\Big(\log \textstyle\sum_c e^{z_{nc}} - z_{n y_n}\Big),
\qquad
\frac{\partial L}{\partial z} = \frac{1}{N}\big(\mathrm{softmax}(z) - \mathrm{onehot}(y)\big).
\end{equation}
%
The gradient reduces to \emph{predicted probability minus true probability}:
all terms of the softmax Jacobian cancel against the logarithm. This is why the
two are always fused. The row maximum is subtracted before exponentiating,
since $e^{z}$ overflows above $z \approx 88$ in single precision.

\begin{lstlisting}[caption={Linear layer and the reverse sweep -- no autograd.}]
class Linear:
    def forward(self, x):
        self.x = x                  # needed for dL/dW
        return x @ self.W + self.b

    def backward(self, g):
        self.dW = self.x.T @ g      # (in,N) @ (N,out) -> (in,out)
        self.db = g.sum(dim=0)      # bias broadcast over the batch
        return g @ self.W.T         # pass gradient down

class MLP:
    def backward(self, g):
        for layer in reversed(self.layers):
            g = layer.backward(g)
\end{lstlisting}

\subsection{Gradient verification}

Both checks report the same quantity,
%
\begin{equation}
\mathrm{rel} = \frac{|\text{analytic} - \text{reference}|}
                    {|\text{analytic}| + |\text{reference}| + 10^{-12}},
\end{equation}
%
so they can be compared directly. Both run in \texttt{float64}. The
finite-difference estimate is centred,
$\partial L/\partial\theta \approx [L(\theta+\varepsilon)-L(\theta-\varepsilon)]/2\varepsilon$
with $\varepsilon = 10^{-5}$, because its error is $\mathcal{O}(\varepsilon^2)$
rather than $\mathcal{O}(\varepsilon)$.

\begin{table}[H]
\centering\small
\caption{Relative error of the hand-written gradients against each reference.}
\begin{tabular}{lccc}
\toprule
Reference method & Worst rel.\ error & Median rel.\ error & Verdict \\
\midrule
Autograd                    & $1.0\times10^{-11}$ & $2.5\times10^{-16}$ & match \\
Centred finite differences   & $3.7\times10^{-9}$  & $1.0\times10^{-10}$ & match \\
\bottomrule
\end{tabular}
\end{table}

\texttt{torch.allclose} returns \texttt{True} for all six gradient tensors. The
worst relative error against autograd exceeds its median only because a few
gradient entries are themselves near zero, where a round-off-sized absolute
difference is a large fraction of a tiny value; the largest \emph{absolute}
difference is $3.3\times10^{-16}$. Finite differences sit near $10^{-10}$
because the subtraction in the numerator cancels most significant digits --
this is a noise floor, not a defect.

Two practical notes. In \texttt{float32} the same check gives relative errors
near $10^{-3}$, which looks like a bug but is only rounding. Also, ReLU is not
differentiable at zero, so a pre-activation lying within $\varepsilon$ of zero
can report a large error even when the code is correct.

\subsection{Training on flattened images}

Each image is flattened to 3072 values and fed to a 3072--512--256--10 network
trained with hand-written SGD.

\fig{f_t1c}{1.0}{Training loss and accuracy. The loss starts near
$\ln 10 = 2.303$, which is what a network that knows nothing must output, and
falls to 0.157. Validation accuracy stalls near 50\% while training accuracy
reaches 94.8\%.}

The network reaches \textbf{50.6\%} best validation accuracy with 1{,}707{,}274
parameters, and overfits heavily: training accuracy ends 46 points higher.

The ceiling matters more than the gap. A fully-connected network cannot get far
past 50--55\% on CIFAR-10 regardless of width or training time, because
flattening destroys locality -- neighbouring pixels become unrelated inputs --
and there is no translation invariance, so a shifted object is an entirely new
input vector. The first layer alone costs $3072 \times 512$ weights and still
sees only one fixed pixel arrangement. This is precisely the gap convolution
closes.

\subsection{Gradient norms at initialisation}

A 10-layer network is given one forward and backward pass at initialisation,
recording $\lVert \partial L/\partial W^{(l)} \rVert$ per layer. Each
activation is paired with the initialiser designed for it: Xavier for sigmoid,
Kaiming for ReLU.

\fig{f_t1d}{0.72}{Per-layer gradient norms, logarithmic axis. Sigmoid
gradients shrink by a factor of about $9.4\times10^{4}$ on the way back to
layer 1; ReLU gradients stay flat.}

Sigmoid collapses because $\sigma'(x) = \sigma(x)(1-\sigma(x)) \le 1/4$, with
equality only at $x=0$. Backpropagation multiplies by one such factor per
layer, so after nine hidden layers the signal is scaled by at most
$(1/4)^9 \approx 4\times10^{-6}$, and in practice far less. The early layers
receive almost no learning signal -- the classic vanishing gradient problem.

ReLU stays flat because its derivative is exactly $0$ or $1$ and never shrinks
what passes through. About half the units are inactive, costing a factor of two
in variance per layer, and Kaiming initialisation restores exactly that factor
by setting $\mathrm{Var}[W] = 2/\text{fan-in}$. The $2$ in the numerator
\emph{is} the ReLU correction, which is why the activation and the initialiser
must be chosen together.

In the ReLU curve layer 1 has the largest norm, but that is a size effect: its
weight matrix is $3072 \times 128$ while the others are $128 \times 128$.

% =====================================================================
\section{Task 2: Optimizers}

\subsection{Implementation and verification}

Both optimizers subclass \texttt{torch.optim.Optimizer}, looping over
\texttt{param\_groups} and keeping per-parameter state in \texttt{self.state}.
Momentum accumulates a velocity, $b_t = \mu b_{t-1} + g_t$, which builds to
$1/(1-\mu) = 10$ along consistent directions while cancelling on directions
that alternate. Adam keeps a running mean and variance and divides one by the
square root of the other:
%
\begin{equation}
m_t = \beta_1 m_{t-1} + (1-\beta_1) g_t, \qquad
v_t = \beta_2 v_{t-1} + (1-\beta_2) g_t^2,
\end{equation}
\begin{equation}
\theta_{t+1} = \theta_t - \eta \,
\frac{m_t / (1-\beta_1^{\,t})}{\sqrt{v_t / (1-\beta_2^{\,t})} + \epsilon}.
\end{equation}
%
Bias correction is necessary because $m$ and $v$ start at zero: at $t=1$,
$m_1 = (1-\beta_1) g_1$ is ten times too small, and dividing by
$1-\beta_1^{\,t}$ restores it exactly.

\begin{lstlisting}[caption={Adam's update, with bias correction.}]
m.mul_(b1).add_(g, alpha=1 - b1)          # first moment
v.mul_(b2).addcmul_(g, g, value=1 - b2)   # second moment
bc1, bc2 = 1 - b1 ** t, 1 - b2 ** t       # bias corrections
denom = (v / bc2).sqrt_().add_(eps)
p.addcdiv_(m / bc1, denom, value=-lr)
\end{lstlisting}

Two details are needed to match PyTorch: the momentum buffer is initialised to
the gradient itself rather than zero, and weight decay is folded into the
gradient as $g \leftarrow g + \lambda\theta$ before the moments are updated.

\begin{table}[H]
\centering\small
\caption{Maximum parameter difference against \texttt{torch.optim} after 50
steps in \texttt{float64}, from identical initial weights and data.}
\begin{tabular}{lcc}
\toprule
Configuration & Max difference & \texttt{allclose} \\
\midrule
SGD (plain)                     & $0.0$               & yes \\
SGD + momentum 0.9              & $0.0$               & yes \\
SGD + momentum + weight decay    & $0.0$               & yes \\
Adam                            & $4.2\times10^{-17}$ & yes \\
Adam, $\beta = (0.8, 0.99)$      & $5.6\times10^{-17}$ & yes \\
Adam + weight decay             & $2.8\times10^{-17}$ & yes \\
\bottomrule
\end{tabular}
\end{table}

The SGD variants are bit-for-bit identical. Adam differs at the seventeenth
digit because PyTorch evaluates $\sqrt{v}/\sqrt{1-\beta_2^t}$ where the
implementation above evaluates $\sqrt{v/(1-\beta_2^t)}$ -- algebraically the
same, but a different order of floating-point operations.

\subsection{Trajectories on an ill-conditioned quadratic}

For $f(x,y) = \tfrac12(x^2 + 50y^2)$ the gradient is $(x, 50y)$, so the surface
is fifty times steeper in $y$. Gradient descent multiplies the $y$ error by
$(1-50\eta)$ each step and therefore only converges while $\eta < 2/50 = 0.04$,
but along $x$ it multiplies by $(1-\eta)$, so progress in the direction that
matters is slow. The steep direction dictates the step size.

\fig{f_t2b}{1.0}{Trajectories from $(-9,1)$ over 60 steps. SGD and momentum use
$\eta = 0.038$, just under the stability limit; Adam uses $\eta = 0.3$ because
its normalised step is roughly $\eta$ regardless of gradient magnitude.}

\begin{table}[H]
\centering\small
\caption{Behaviour on the narrow valley.}
\begin{tabular}{lcc}
\toprule
Optimizer & Sign changes in $y$ & Final $f$ \\
\midrule
SGD              & 60 of 60 & 0.388 \\
SGD + momentum   & 30       & 0.058 \\
Adam             & 6        & 0.109 \\
\bottomrule
\end{tabular}
\end{table}

SGD changes direction on every step: with $\eta = 0.038$ the $y$ error is
multiplied by $-0.9$, flipping sign while shrinking only 10\%, and $x$ crawls
from $-9$ to $-0.88$. Momentum halves the oscillation and travels five times
further at the same learning rate, because alternating components cancel in the
velocity buffer while consistent ones accumulate. Adam barely oscillates, since
dividing by $\sqrt{\hat v}$ normalises each coordinate by its own gradient
history and makes the conditioning largely irrelevant.

Momentum reaches the lowest final value. On a clean convex problem, well-tuned
momentum is hard to beat; Adam's advantage appears when gradients are noisy and
badly scaled, as in the next section.

\subsection{Learning-rate sensitivity}

A small CNN with batch norm (141{,}034 parameters) is trained for 15 epochs with
each optimizer at the same three learning rates, so that the position and
sharpness of each peak can be read off directly.

\fig{f_t2c}{0.72}{Best validation accuracy against learning rate. Adam peaks
inside the grid; both SGD variants are still rising at the right edge.}

\begin{table}[H]
\centering\small
\caption{Best validation accuracy (\%) after 15 epochs.}
\begin{tabular}{lcccc}
\toprule
Optimizer & $\eta = 0.001$ & $\eta = 0.01$ & $\eta = 0.1$ & Spread \\
\midrule
SGD            & 42.76 & 62.54 & \textbf{71.58} & \textbf{28.8} \\
SGD + momentum & 64.22 & 75.28 & \textbf{79.14} & 14.9 \\
Adam           & 76.78 & \textbf{80.40} & 69.14 & 11.3 \\
\bottomrule
\end{tabular}
\end{table}

Sensitivity is the main result. Plain SGD swings 28.8 points across the grid and
is nearly useless at $\eta = 0.001$, while Adam varies by only 3.6 points
between $\eta = 0.001$ and $\eta = 0.01$. A careless learning rate costs almost
nothing with Adam and close to thirty points with plain SGD, which is why Adam
is the usual default when there is no budget for tuning.

Adam is nevertheless the only optimizer that fails outright, collapsing by 11
points at $\eta = 0.1$: its step is about $\eta$ per parameter irrespective of
gradient size, so updates become as large as the weights themselves.

Both SGD curves are still rising at $\eta = 0.1$, so this grid does not bracket
their optima and 71.58\% should be read as a lower bound. The tenfold
amplification from momentum is visible at matched quality instead: momentum at
$\eta = 0.01$ reaches 75.28\%, beating plain SGD at $\eta = 0.1$ (71.58\%) with
one tenth of the learning rate.

Finally, this CNN reaches 80.40\% with 141{,}034 parameters where the Task 1 MLP
reached 50.6\% with 1{,}707{,}274 -- twelve times fewer parameters for thirty
points more accuracy.

% =====================================================================
\section{Task 3: Regularization and residual connections}

\subsection{Regularization on 5{,}000 images}

The training set is reduced to 5{,}000 images, because with all 45{,}000 a
141k-parameter network barely overfits and the effect would be too small to
measure. Cutout zeroes a random $8\times8$ square after normalisation, so the
hole takes the mean colour. Mixup blends two images and their labels in the
same proportion,
%
\begin{equation}
\tilde x = \lambda x_i + (1-\lambda) x_j, \qquad
L = \lambda\,\mathrm{CE}(f(\tilde x), y_i) + (1-\lambda)\,\mathrm{CE}(f(\tilde x), y_j),
\end{equation}
%
with $\lambda \sim \mathrm{Beta}(1,1)$, forcing linear behaviour between
training examples rather than overconfident predictions in the gaps.

\fig{f_augment}{1.0}{Each augmentation applied to one image. Flip and crop
produce images that are still clearly the same object; cutout hides evidence;
mixup produces an image that does not exist, labelled in two classes.}

\begin{lstlisting}[caption={Cutout and mixup, batched on the GPU.}]
def apply_cutout(x, length, gen):
    cy = torch.randint(0, 32, (x.shape[0],), device=DEVICE, generator=gen)
    cx = torch.randint(0, 32, (x.shape[0],), device=DEVICE, generator=gen)
    half = length // 2
    my = (ARANGE32[None,:] >= cy[:,None]-half) & (ARANGE32[None,:] < cy[:,None]+half)
    mx = (ARANGE32[None,:] >= cx[:,None]-half) & (ARANGE32[None,:] < cx[:,None]+half)
    hole = my[:,:,None] & mx[:,None,:]
    return x * (~hole)[:,None,:,:]

def mixup_batch(x, y, alpha, gen, rng):
    lam  = float(rng.beta(alpha, alpha))
    perm = torch.randperm(x.shape[0], device=DEVICE, generator=gen)
    return lam * x + (1 - lam) * x[perm], y, y[perm], lam
\end{lstlisting}

Training accuracy is always measured on clean, unaugmented images in a separate
pass. Measuring it on augmented batches would make mixup's figure a different
quantity and the gap column meaningless.

\begin{table}[H]
\centering\small
\caption{Seven configurations after 80 epochs on 5{,}000 images.}
\begin{tabular}{lcccc}
\toprule
Configuration & Train (\%) & Val (\%) & Gap & Best val (\%) \\
\midrule
Baseline      & 100.00 & 67.36 & 32.6 & 67.36 \\
Weight decay  & 100.00 & 68.08 & 31.9 & 68.42 \\
Dropout       & 100.00 & 67.44 & 32.6 & 67.80 \\
Flip + crop   & 83.10  & 65.72 & 17.4 & \textbf{70.88} \\
Cutout        & 98.06  & 63.94 & 34.1 & 67.94 \\
Mixup         & 99.46  & 62.94 & 36.5 & 65.46 \\
All combined  & 72.56  & 65.02 & \textbf{7.5} & 68.82 \\
\bottomrule
\end{tabular}
\end{table}

\fig{f_t3a}{1.0}{Train versus validation accuracy, and the generalisation gap.}

The training column gives the clearest result: baseline, weight decay and
dropout all reach exactly 100.00\%, memorising every image despite the
regularization. Only the methods that change the data prevent this. On this
dataset, constraining the weights barely helped while changing the images did.

Dropout contributed almost nothing because it sits only after global average
pooling, on a 128-dimensional vector -- far weaker than dropout between
convolutional layers. That is a consequence of placement, not of the method.

Two results deserve comment. Mixup \emph{reduced} accuracy by 1.9 points: with
$\alpha = 1$ on only 5{,}000 images nearly every training image is a blend and
the network sees little clean signal. And ``all combined'' is not the best
configuration -- it has by far the smallest gap (7.5 points) but only 68.82\%,
below flip and crop alone at 70.88\%, with training accuracy down to 72.56\%.
Five regularizers over-constrain a small network, which is underfitting rather
than successful regularization. Closing the gap is not the same as generalising
better.

\subsection{Plain versus residual networks}

Networks follow the CIFAR ResNet form: a stem convolution, three stages of $n$
blocks at 16/32/64 channels, global average pooling and a linear head, giving
depth $6n+2$. Shortcuts use Option A -- stride subsampling with zero-padded
channels -- which has no parameters, so the two families have identical
parameter counts and the only difference is one addition.

\begin{lstlisting}[caption={The only difference between the two families.}]
def forward(self, x):
    out = F.relu(self.b1(self.c1(x)))
    out = self.b2(self.c2(out))
    if self.residual:
        out = out + self.shortcut(x)   # <-- the only difference
    return F.relu(out)
\end{lstlisting}

No augmentation is used here. The question is whether a deep network can be
optimised, so the networks are allowed to fit the training data as hard as they
can and the training losses are compared directly.

\begin{table}[H]
\centering\small
\caption{Final training loss after 30 epochs. Parameter counts are identical
within each depth.}
\begin{tabular}{lccc}
\toprule
Network & Parameters & Final train loss & Val (\%) \\
\midrule
Plain-8     & 75{,}290  & 0.0617 & 79.40 \\
Residual-8  & 75{,}290  & 0.0641 & 79.10 \\
Plain-20    & 269{,}722 & 0.0037 & 82.56 \\
Residual-20 & 269{,}722 & 0.0017 & 83.16 \\
Plain-56    & 853{,}018 & \textbf{0.7719} & 57.72 \\
Residual-56 & 853{,}018 & \textbf{0.0009} & 80.78 \\
\bottomrule
\end{tabular}
\end{table}

\fig{f_t3b}{1.0}{Training loss per epoch, logarithmic axis. The 56-layer plain
network gets stuck well above its own 20-layer version.}

\fig{f_t3b_summary}{0.62}{Final training loss against depth. Deeper always
helps the residual family; for the plain family, 20 to 56 layers makes training
loss 200 times worse.}

The plain network breaks at depth 56: its training loss rises from 0.0037 to
0.7719, roughly 200 times worse, despite having three times more parameters.
The residual network at the same depth and parameter count reaches 0.0009, a
factor of about 860 better.

This is not overfitting. Overfitting means low training loss with poor
validation accuracy, whereas here the training loss itself is bad, on data the
network is being fitted to. A 56-layer network can represent anything a
20-layer one can, by setting the extra layers to the identity, so it has
strictly more capacity and must be able to match plain-20. It fails to, so the
capacity exists and gradient descent cannot reach it. This is the degradation
problem.

A residual block computes $x + F(x)$, so representing the identity only requires
driving $F$ to zero, which weight decay does readily. A plain block must learn
to reproduce its input before improving on it. The shortcut also gives the
gradient a path of derivative exactly one back to earlier layers -- the clean
route that Task 1(d) showed sigmoid lacked.

At depth 8 the two families are effectively tied, with plain marginally ahead.
That confirms the benefit is specifically about depth and is not free at any
size. Residual-56 also reached the lowest training loss of all six networks but
not the best validation accuracy, which belongs to residual-20; with 30 epochs
and no augmentation the deepest network fits hardest and generalises slightly
worse.

% =====================================================================
\section{Task 4: Best CNN under 500k parameters}

\subsection{Design}

Twenty epochs is the binding constraint. Task 3(b) showed that at 30 epochs a
56-layer network had still not overtaken a 20-layer one on validation accuracy,
so the design starts from depth 20 and spends the remaining budget on width.

\begin{table}[H]
\centering\small
\caption{Choosing the widest configuration that fits the budget.}
\begin{tabular}{lccc}
\toprule
Channels & Depth & Parameters & Fits 500k \\
\midrule
16 / 32 / 64          & 20 & 269{,}722 & yes, 46\% unused \\
\textbf{21 / 42 / 84} & \textbf{20} & \textbf{463{,}627} & \textbf{yes, 92.7\% used} \\
22 / 44 / 88          & 20 & 508{,}672 & no \\
24 / 48 / 96          & 14 & 386{,}626 & yes, but shallower \\
20 / 40 / 80          & 26 & 572{,}430 & no \\
\bottomrule
\end{tabular}
\end{table}

The final recipe is a residual network of depth 20 with widths 21/42/84, trained
with SGD and momentum 0.9, learning rate 0.1 annealed to zero by a cosine
schedule, weight decay $5\times10^{-4}$, and flip, crop and cutout augmentation,
at batch size 128 for 20 epochs. Every ingredient except the schedule comes from
an earlier task; cosine annealing is used because a 20-epoch budget leaves no
room for a step schedule's plateaus.

\subsection{Ablation}

Each row removes exactly one ingredient and changes nothing else.

\begin{table}[H]
\centering\small
\caption{Ablation of four design choices, ranked by cost of removal.}
\begin{tabular}{lccc}
\toprule
Configuration & Best val (\%) & Change & Final train loss \\
\midrule
Full recipe          & \textbf{89.60} & ---              & 0.2281 \\
Minus augmentation   & 85.48 & $\mathbf{-4.12}$ & \textbf{0.0047} \\
Adam instead of SGD  & 86.90 & $-2.70$          & 0.3284 \\
Minus shortcut       & 87.56 & $-2.04$          & 0.3228 \\
Minus weight decay   & 87.58 & $-2.02$          & 0.2848 \\
\bottomrule
\end{tabular}
\end{table}

\fig{f_t4}{1.0}{What each ingredient contributes, and training loss per epoch.
The curve for ``minus augmentation'' falls far below the others while scoring
worst on validation.}

Augmentation matters roughly twice as much as anything else, and the training
loss explains why: without it the loss falls to 0.0047, about 48 times lower
than the full recipe, while validation accuracy drops four points. Fitting the
training data much better and the held-out data much worse is the signature of
overfitting.

Removing the shortcut works through a different mechanism. Training loss rises
to 0.3228, so the plain network fits the training data \emph{less} well -- an
optimisation cost rather than a generalisation one, consistent with Task 3(b).
Both variants have identical parameter counts. The effect is smaller than at
depth 56 because depth 20 is comfortably trainable without shortcuts.

Shortcut and weight decay differ by only 0.02 points and should be treated as
tied rather than ranked, since that is well below run-to-run variation from a
single seed. SGD beat Adam here, reversing Task 2(c): Adam's strength is
tolerance of a poorly chosen learning rate, but once the rate is well chosen and
annealed, tuned momentum overtakes it, and Adam also ends at a higher training
loss.

\subsection{Final result}

\begin{table}[H]
\centering\small
\caption{Final model, with the test set evaluated once.}
\begin{tabular}{ll}
\toprule
Architecture & Residual, depth 20, widths 21/42/84 \\
Parameters   & 463{,}627 of 500{,}000 allowed (92.7\%) \\
Epochs       & 20 of 20 allowed \\
Validation accuracy & 89.60\% \\
\textbf{Test accuracy} & \textbf{90.05\%} \\
\bottomrule
\end{tabular}
\end{table}

Test accuracy came in 0.45 points above validation. The validation split
absorbed every design decision and therefore carries some optimistic bias, but
with only five configurations compared that bias is evidently smaller than the
difference between the two splits; the test set is also twice the size and so
the less noisy measurement. The evaluation is cached, so re-running the notebook
reads the stored value rather than measuring again -- once a test number
influences a decision it stops being a held-out estimate.

The binding constraint was epochs, not parameters: the final training loss is
0.2281 and the validation curve is still rising at epoch 20, so the model is
underfitted. With more budget the order worth trying is longer training first,
then more width, and more depth last.

% =====================================================================
\section{Conclusion}

\begin{table}[H]
\centering\small
\caption{Progression across the four tasks.}
\begin{tabular}{llcc}
\toprule
Task & Model & Parameters & Accuracy \\
\midrule
1(c) & MLP on flattened pixels     & 1{,}707{,}274 & 50.6\% val \\
2(c) & Small CNN, tuned optimizer   & 141{,}034     & 80.4\% val \\
3(b) & Residual-20, no augmentation & 269{,}722     & 83.2\% val \\
4    & Residual net, full recipe    & 463{,}627     & \textbf{90.1\% test} \\
\bottomrule
\end{tabular}
\end{table}

Task 1 showed that a fully-connected network cannot exceed roughly half on
images, because flattening discards locality and translation invariance. Task 2
showed that the choice of optimizer and learning rate is worth tens of points on
the same architecture, and that Adam's main practical value is insensitivity to
that choice. Task 3 separated two distinct failures -- memorisation from too
little data, fixed by changing the data, and untrainability from too many
layers, fixed by shortcuts. Task 4 combined these under a fixed budget.

The recurring theme is that structure matters more than size. The final model
has 3.7 times fewer parameters than the Task 1 MLP and is 39 points more
accurate, because convolution, residual connections and augmentation each encode
something true about images that the first network had to learn from data and
could not.

% =====================================================================
\section*{Acknowledgement of AI tool use}

Claude (Anthropic) was used as a coding assistant while preparing the notebook:
scaffolding the layer classes, the gradient-checking harness and the plotting
code. All implementations were reviewed, executed and verified by me, and the
interpretation of the results is my own. I understand every line of the
submitted code.

\end{document}
